% !TEX root = ../../../main.tex
% Block card: detection.sax_csax -- "Symbolic encoding (cSAX)" (Part I, the blocks chapter).
% This card carries the plain explanation of the whole SAX family; the pSAX and dSAX cards
% point back to it.
\begin{plainterms}[Turning a signal into letters (SAX)]
Imagine summarising a long walk by writing one letter per minute: \emph{a} for low ground,
\emph{b} for middling, \emph{c} for high. That is what the SAX family does to a recording.
First it cuts the span into equal pieces and replaces each piece by its average (a
\emph{piecewise average}). Then it draws a few horizontal lines, the \emph{breakpoints}, and
gives each average the letter of the band it falls in. In classic SAX
\citep{lin2003symbolic} the breakpoints are placed so that, if the signal were bell-shaped
(normally distributed), every letter would turn up equally often. Once a recording is a
string, you can search it like text. The three blocks here differ in how they draw the
lines: \textbf{cSAX} finds the levels the averages cluster around, \textbf{pSAX} fits a
fixed number of levels to the averages' spread, and \textbf{dSAX} codes each piece's
\emph{rise} instead of its average.
\end{plainterms}

\cardfield{What it does}
Turns a span into a string of letters, one per piece of time, with the letters set by the
levels the signal's piece-averages naturally cluster around.

\cardfield{When to use it}
When the signal sits at a few distinct levels (two resting potentials, say) and you want the
string to name which level each stretch was at, without choosing how many letters
beforehand. For shapes (rises and falls) use \blockref{detection.sax_dsax}; for a fixed
number of letters use \blockref{detection.sax_psax}.

\cardfield{What goes in, what comes out}
In goes a span of one channel. Out comes a string: one letter per piece, \emph{a} for the
lowest level and onwards. \blockref{detection.symbol_search} can then search it.

\cardfield{Settings that matter}
\emph{segment\_mode} picks which one control sets the piece length; the default is
\emph{seconds\_per\_symbol}, 20\,s per letter. Beside the settings, the block page shows
live read-outs before anything runs: the span length, the samples and seconds per letter,
how many letters you will get (amber under 10 or over 2000, red under 2) and how many samples
at the end are dropped because they do not fill a whole piece. The block also carries a rule
for span-aware defaults (about one letter per 20 samples, 16 to 256 letters), but the app
does not apply it yet: the page says the hook is declared and not computed.

\cardfield{Watch out for}
The number of letters is not yours to choose, and it is fragile: if the clustering finds
fewer than two levels, the block quietly falls back to classic SAX with ten letters, and the
app does not say so beyond the alphabet size of 10 in the result's summary. The clustering
starts from random guesses; the app always uses the same fixed seed, so a run repeats, but
on one test channel other seeds gave anywhere from 3 to 6 letters. A single
missing sample (NaN) in the span spoils the whole encoding.
